\clearpage
\begin{proofbox}
  \textbf{\textcolor{CProof}{先说明要证明什么}}
  我们要证明的是：\textbf{只要定义域左右成对}，
  “沿 $x=a$ 折叠后图象重合”和“所有成对位置函数值相同”其实是一回事。
  下面按理解难度从低到高给出三种讲法。

  \medskip
  \textbf{\textcolor{CProof}{方法一：折纸法——只用初中轴对称（门槛最低）}}
\begin{center}\begin{tikzpicture}[font=\small,>=stealth,line join=round]
\draw[gray!40,rounded corners=3pt] (0,0) rectangle (14.2,8.3);\node at (7.1,7.700000000000001) {左右等距、函数值同高：对称轴就在中间};
\begin{scope}[shift={(1.7923076923076926,2.038732394366197)},x=1.153846153846154cm,y=0.647887323943662cm]\draw[->,gray] (-0.6,0)--(4.6,0) node[right]{$x$};\draw[->,gray] (0,-0.6)--(0,6.5) node[above]{$y$};\draw[blue,thick,solid,domain=-0.1:4.1,samples=180] plot(\x,{(\x-2)^2+1});\fill[orange!85!black] (0,5) circle(2.2pt);\fill[orange!85!black] (1,2) circle(2.2pt);\fill[orange!85!black] (2,1) circle(2.2pt);\fill[orange!85!black] (3,2) circle(2.2pt);\fill[orange!85!black] (4,5) circle(2.2pt);\draw[violet,dashed] (2,-.5)--(2,6);\node[above,violet] at (2,6){$x=2$};\draw[orange!85!black,dashed] (0,5)--(4,5);\draw[orange!85!black,dashed] (1,2)--(3,2);\foreach \v in {0,1,2,3,4}{\draw[gray] (\v,0)--(\v,-.1);\node[below=5pt] at (\v,-.1){$\v$};}\end{scope}\node[anchor=west,align=left,text width=5.0cm] at (8.4,4.0) {$f=(x-2)^2+1$\\[9pt]$f(0)=f(4)=5$\\[9pt]$f(1)=f(3)=2$\\[9pt]每对横坐标之和为 $4$\\[9pt]轴的位置：$4\div2=2$};
\node[text width=12.6cm,align=center] at (7.1,.65) {图示具体配对；完整证明要求定义域左右成对，且每一对都同高。};
\end{tikzpicture}\end{center}


  \begin{enumerate}[leftmargin=*,itemsep=0.5em]
    \item 取图象上一点 $P$，让它的横坐标为 $a+t$，那么
      $P=(a+t,f(a+t))$。
    \item 沿着竖直线 $x=a$ 对折。点 $P$ 离竖线水平距离为 $|t|$，
      对折后在另一侧同样远的位置，横坐标变成 $a-t$，高度却不改变。
      因此折过去的点是 $(a-t,f(a+t))$。
    \item \textbf{如果原图确实左右重合}，折过去的点也一定在原图上。
      而原图在横坐标 $a-t$ 处的高度叫做 $f(a-t)$，于是
      \[
        f(a+t)=f(a-t).
      \]
    \item \textbf{反过来}，如果定义域左右成对，并且每一对位置的高度都相同，
      那么图象上任意一点折过去，都能找到原来图象上的对应点。
      所有点都能对应，整幅图象便关于 $x=a$ 对称。
  \end{enumerate}
  \textbf{这一步最关键：}不能只看两三对点；这里的“任意一点”意味着整条图象。

  \medskip
  
\end{proofbox}
\clearpage
\begin{proofbox}
\textbf{\textcolor{CProof}{方法二：中点法——两边横坐标的平均数是 $a$（做题最快）}}
\begin{center}\begin{tikzpicture}[font=\small,>=stealth,line join=round]
\draw[gray!40,rounded corners=3pt] (0,0) rectangle (14.2,8.5);\node at (7.1,7.9) {例6：真实代入位置相加恒为6，中点才是轴};
\draw[->,gray] (1,2)--(13,2) node[right]{$x$};\draw[violet,dashed] (7,1.5)--(7,6.7);\draw[blue,thick] (3,4)--(11,4);\fill[orange!85!black] (3,4) circle(2.4pt);\draw[dashed,gray] (3,2)--(3,4);\node[below=7pt] at (3,2){$t+1$};\fill[orange!85!black] (11,4) circle(2.4pt);\draw[dashed,gray] (11,2)--(11,4);\node[below=7pt] at (11,2){$5-t$};\node[above,violet] at (7,6.7){$x=3$};\node[above=8pt] at (10,4){同高};
\node[text width=12.6cm,align=center] at (7.1,.65) {横向以 $t=0$ 的配对 $1,5$ 示意；一般位置之和恒为6，不指定函数值。};
\end{tikzpicture}\end{center}


  \begin{enumerate}[leftmargin=*,itemsep=0.5em]
    \item 在图象上取点 $P=(x,f(x))$。它在镜子另一侧的点的横坐标设为 $u$。
    \item 这两个横坐标中间恰好是 $a$，所以
      \[
        \frac{x+u}{2}=a \quad\Longrightarrow\quad u=2a-x.
      \]
      对折不会改变高度，故镜像位置应为 $(2a-x,f(x))$。
    \item 如果图象关于 $x=a$ 对称，该镜像点在图象上，
      因此 $f(2a-x)=f(x)$。
    \item 反过来，若\textbf{所有} $x$ 都有 $2a-x$ 也在定义域中，且
      $f(2a-x)=f(x)$，则任取 $P$，它的镜像点恰好也在图象上。
      所以图象可以沿 $x=a$ 折叠重合。
  \end{enumerate}
  \textbf{得到同一个结论：}
  \[
    \boxed{f(x)=f(2a-x)}.
  \]
  把 $x$ 换成 $a+t$，右边就变成 $a-t$，于是得到 $f(a+t)=f(a-t)$。
  \textbf{用法：}知道一侧的值时，直接把自变量换成 $2a-x$，不用重新算。

  \medskip
  
\end{proofbox}
\clearpage
\begin{proofbox}
\textbf{\textcolor{CProof}{方法三：平移法——把对称轴移到 $0$，联系偶函数}}
\begin{center}\begin{tikzpicture}[font=\small,>=stealth,line join=round]
\draw[gray!40,rounded corners=3pt] (0,0) rectangle (14.2,8.5);\node at (7.1,7.9) {平移证明：把原来的对称轴搬到零点};
\node at (3.2,6.9){$f(x)=(x-2)^2+1$};\begin{scope}[shift={(1.2,2.4)},x=1cm,y=.62cm]\draw[->,gray] (-0.5,0)--(4.5,0);\draw[->,gray] (0,0)--(0,6);\draw[blue,thick,domain=0:4,samples=120] plot(\x,{(\x-2)^2+1});\draw[violet,dashed] (2,0)--(2,5.7);\node[below=6pt] at (2,0){$2$};\fill[orange!85!black] (2,1) circle(2pt);\end{scope}\node at (10.3,6.9){$g(t)=f(2+t)=t^2+1$};\begin{scope}[shift={(10.3,2.4)},x=1cm,y=.62cm]\draw[->,gray] (-2.5,0)--(2.5,0);\draw[->,gray] (0,0)--(0,6);\draw[blue,thick,domain=-2:2,samples=120] plot(\x,{(\x)^2+1});\draw[violet,dashed] (0,0)--(0,5.7);\node[below=6pt] at (0,0){$0$};\fill[orange!85!black] (0,1) circle(2pt);\end{scope}\draw[->,orange!85!black,thick] (6.1,4.3)--(7.2,4.3);\node at (7.1,1.5){原图向左平移2个单位；两图横纵单位分别相同，曲线高度不变。};
\node[text width=12.6cm,align=center] at (7.1,.65) {横坐标由 $x$ 改为 $t=x-2$，因此轴 $x=2$ 对应新的 $t=0$。};
\end{tikzpicture}\end{center}


  \begin{enumerate}[leftmargin=*,itemsep=0.5em]
    \item 假设原图象想要以 $x=a$ 为对称轴。我们改用 $t=x-a$ 表示横坐标，
      等于把“原来的 $a$ 位置”当成新的 $0$。
    \item 记 $g(t)=f(a+t)$。那么新图象在 $t$ 和 $-t$ 处的高度分别是
      \[
        g(t)=f(a+t),\qquad g(-t)=f(a-t).
      \]
    \item 如果 $f(a+t)=f(a-t)$ 对所有允许的 $t$ 成立，
      那么 $g(t)=g(-t)$。这正是高中学过的偶函数的定义：
      新图象关于 $t=0$ 对称。
    \item 把横坐标原点再移回去，$t=0$ 对应 $x=a$，
      所以原图象关于 $x=a$ 对称。反过来按同样步骤倒着做，也成立。
  \end{enumerate}
  \textbf{注意：}平移前后都要保留完整的定义域。原定义域关于 $a$ 对称，
  新定义域才关于 $0$ 对称。

  \medskip
  \textbf{\textcolor{CProof}{三种方法怎么选？}}
  \begin{itemize}[leftmargin=*,itemsep=0.35em]
    \item 想让初中生一听就懂：选\textbf{折纸法}。
    \item 已知对称轴，要求值、求未知数或求另一半解析式：选\textbf{中点法}。
    \item 题目同时出现偶函数、平移后的函数：选\textbf{平移法}。
  \end{itemize}
\end{proofbox}
